\documentclass[11pt]{article}
\usepackage{mathblatt}
% Präambel des Setzers (werkzeuge/setzer-praeambel.tex, wortgleich); dazu \gteil (Muster T2U)
\makeatletter
\setlength{\mbpfbe}{0mm}% keine Punkte (pruefung.md § 5: nur beim Prüfstein)
\newlength{\gEin}\setlength{\gEin}{0pt}
\let\gSchrift\relax
\renewcommand{\pfkreuzzeile}[1]{\par\noindent\hangindent3.2em\hangafter1\hspace*{1.6em}$\square$\ #1\par}
\newcommand{\gkreuz}[1]{$\square$\ #1\hspace{2em}}
\newcommand{\gkreuzzeile}[2]{\par\noindent\hangindent2.2em\hangafter1\mbox{#1\,$\square$}\ #2\par}% Buchstabe, Option: eine je Zeile (neben Skizze)
% links, was man ansieht, rechts, was man tut; Spaltengrenze überall an derselben Stelle (gemeinsam.md „Satz“)
% v0.17: Fußeintrag „Nr: Wert“ – kurz in einer Box (nie getrennt), lang umbrechbar mit der Nummer am Anfang
% (vorher \mbox: lange Lösungen liefen über den Rand, AS40 138 pt); robust, weil \fusshilfe in die .aux schreibt
\NewDocumentCommand{\gfuss}{m m}{\sbox0{#1:~#2}\ifdim\wd0<0.45\textwidth\mbox{#1:~#2}\else\mbox{#1:}\nobreak\ #2\fi}
\newsavebox{\gbox}\newlength{\gLinks}
% v0.10 (Lehrer 13.10.): keine starre Halbierung – die linke Spalte ist so breit wie die Figur (+ 6 mm), die
% Antwort steht direkt daneben; ist die Figur breiter als 0,62 Zeile, steht die Antwort darunter
\newcommand{\gzwei}[2]{\par\smallskip\noindent\sbox{\gbox}{#1}%
  \setlength{\gLinks}{\dimexpr\wd\gbox+6mm\relax}%
  \ifdim\gLinks>0.62\linewidth
    \usebox{\gbox}\par\smallskip\noindent\begin{minipage}{\linewidth}#2\end{minipage}%
  \else
    \begin{minipage}[t]{\gLinks}\vspace{0pt}\usebox{\gbox}\end{minipage}%
    \begin{minipage}[t]{\dimexpr\linewidth-\gLinks\relax}\vspace{0pt}\raggedright #2\end{minipage}%
  \fi\par}
% v0.9 (12.10., pruefung.md § 5 „Marke mit B, Bedingungen rechts“): drittes Argument von pfaufg = Bedingungen
% der Prüfung (v0.10: \pfohnetr = TR durchgestrichen, „FOR“), klein grau im rechten Rand auf der Grundlinie der Nummer, Zeilen mit \\
\newcommand{\gbedrand}[2]{\ifblank{#2}{}{\rlap{\hspace*{\dimexpr#1+2mm\relax}%
  \parbox[t]{15mm}{\raggedright\scriptsize\leavevmode\color{mbgrau}\strut#2}}}}
% Teil eines Bündels (a, b …): Marke links im Rand, Bedingungen rechts im Rand, auf der Zeile des Teils
\newcommand{\gteilrand}[2]{\llap{\mbpfrandmarke{#1}\hspace*{\mbpfnr}}\gbedrand{\linewidth}{#2}}
% Teil eines Bündels: Buchstabe in eigener Spalte, Inhalt eingerückt (\gEin = Einzug, damit die Spaltengrenze von
% \gzwei an derselben Stelle bleibt)
\newenvironment{gbteil}[3]{\par\smallskip\noindent\gteilrand{#2}{#3}\makebox[\mbpfteil][l]{\textbf{#1}}%
  \begin{minipage}[t]{\dimexpr\linewidth-\mbpfteil\relax}\setlength{\gEin}{\mbpfteil}\raggedright\strut\ignorespaces}%
 {\end{minipage}\par}
% Grafik höchstens so breit wie die Spalte (Bündel nebeneinander, Koordinatensystem ohne Karo)
\newsavebox{\gpbox}
\newcommand{\gpassend}[1]{\sbox{\gpbox}{#1}\ifdim\wd\gpbox>\linewidth\typeout{GPASSEND: verkleinert \the\wd\gpbox\space auf \the\linewidth}\resizebox{\linewidth}{!}{\usebox{\gpbox}}\else\usebox{\gpbox}\fi}
% Skizzenfeld (v0.10, Lehrer 13.10.): Karo ab dem linken Rand über die Zeile, „Skizze“ klein grau oben links im Feld
% (wie \pfskizzenfeld der Vorlage); der rechte Rand bleibt den Bedingungen
\newcommand{\gskizze}[1]{\par\smallskip\noindent\begin{tikzpicture}
  \clip (0,0) rectangle ({\linewidth-1mm},{#1*10mm});
  \draw[black!16,line width=0.3pt,step=5mm] (0,0) grid ({\linewidth},{#1*10mm});
  \node[anchor=north west,inner sep=2pt,font=\scriptsize,text=mbgrau,fill=white] at (0,{#1*10mm}) {Skizze};
  \end{tikzpicture}\par}
\RenewDocumentEnvironment{pfaufg}{m m m}{%
  \begin{lrbox}{\mbaufgabebox}\begin{minipage}[t]{\linewidth}%
  \gSchrift\noindent\hspace*{\gEin}\mbpfrandmarke{#1}\gbedrand{\linewidth-\gEin}{#3}%
  \begin{minipage}[t]{\mbpfnr}\strut\textbf{#2}\end{minipage}%
  \begin{minipage}[t]{\dimexpr\linewidth-\gEin-\mbpfnr\relax}\raggedright\strut\ignorespaces}%
 {\end{minipage}%
  \end{minipage}\end{lrbox}%
  \setlength{\mbaufgabehoehe}{\dimexpr\ht\mbaufgabebox+\dp\mbaufgabebox\relax}%
  \ifdefempty{\gkopfnach}{\par\addvspace{5pt}\Needspace*{\mbaufgabehoehe}}{%
    % Stufenkopf, der mit dieser (ersten) Aufgabe zusammenbleibt: Platz für beide verlangen, dann den Kopf setzen
    % (12.10., Quadratische: ein Kopf stand allein unten, die Aufgabe mit Koordinatensystem auf der nächsten Seite)
    % v0.17 (AS40 „Terme deuten“ allein unten): hinter dem Kopf steht das Sprungziel (\special, ein Knoten, nach
    % dem die Seite brechen durfte) – darum \nobreak nach dem Kopf; Vorhalt 56 pt = Kopf mit Linie und Abständen
    \par\Needspace*{\dimexpr\mbaufgabehoehe+56pt\relax}\gkopfnach\par\nobreak\gdef\gkopfnach{}\addvspace{5pt}}%
  \noindent\usebox{\mbaufgabebox}\par\addvspace{5pt}}
% Stufenkopf vormerken: \gkopfvor{<Kopf>} – die nächste pfaufg setzt ihn, sobald beide zusammen Platz haben
\def\gkopfnach{}
\newcommand{\gkopfvor}[1]{\gdef\gkopfnach{#1}}
% Teilaufgabe mit Skizze links in gzwei: Buchstabe in eigener Spalte, darunter Text und Skizze
\newcommand{\gteil}[2]{\makebox[\mbpfteil][l]{\textbf{#1}}\parbox[t]{\dimexpr0.5\textwidth-\gEin-\mbpfnr-5mm-\mbpfteil\relax}{\raggedright #2}}
% Stufenkopf der Vorlage (mathblatt.sty 2026-10-12: Linie, Abstand, 13 pt, \pfjahre) ohne ihren Vorhalt von
% 14 Zeilen: der Setzer fordert vor dem Kopf 0,16 Seitenhöhe (Kopf + erste Aufgabe), mehr nicht
\setlength{\mbpfstufevorhalt}{0pt}
\makeatother
\newcommand{\karo}[1]{\par\smallskip\noindent\begin{tikzpicture}
  \clip (0,0) rectangle ({\linewidth-1mm},{#1*10mm});
  \draw[black!16,line width=0.3pt,step=5mm] (0,0) grid ({\linewidth},{#1*10mm});
  \end{tikzpicture}\par}
\newcommand{\kk}[1]{$\square$\,#1\hspace{0.9em}}
\newcommand{\linie}[1]{\,{\color{black!45}\rule[-1pt]{#1}{0.4pt}}\,}
\newcommand{\pkt}{\textperiodcentered{}}
% v0.14 (14.10.): Stufenkopf mit fünf Punkten (Option = Jahre für die Bestellkarte) oder „zuletzt 19“
\NewDocumentCommand{\gpunkte}{o m}{#2}
\newcommand{\gzuletzt}[1]{zuletzt #1}
% v0.7 (12.10.): Schülerfrage unter dem Stufenkopf; Ankreuzoption als Bild (Buchstabe und Kästchen darunter)
\newcommand{\gfrage}[1]{\par\nobreak\noindent\hspace*{\mbpfnr}{\small\itshape #1}\par\nobreak\addvspace{2pt}}
\newlength{\gbildbreite}\setlength{\gbildbreite}{0.32\linewidth}% Setzer setzt sie je Reihe (vier Bilder: 0,24)
\newcommand{\gbildoption}[3]{\begin{minipage}[t][#3][b]{\gbildbreite}\centering #2\par\smallskip\mbox{#1\,$\square$}\end{minipage}}
% Lineare (12.10.): Bild mit Buchstabe ohne Kästchen (Diagramm zuordnen); Aussage mit „wahr / falsch“ zum Ankreuzen
\newcommand{\gbildname}[3]{\begin{minipage}[t][#3][b]{\gbildbreite}\centering #2\par\smallskip\mbox{\textbf{#1}}\end{minipage}}
\newcommand{\gwahrfalsch}[2]{\par\noindent\hangindent2.2em\hangafter1\makebox[2.2em][l]{(#1)}#2\hfill\mbox{wahr $\square$\quad falsch $\square$}\par}
\begin{document}

\pfheftstil
\pflaufkopf[Abitur GK \pkt{} Prüfung]{Binomialwahrscheinlichkeit}{AS70}\pfseitevon
\pfheftkopf{Binomialwahrscheinlichkeit}{}
\pfstufenzeile{Wahrscheinlichkeit gesucht \pkt{} n, p oder k gesucht}

% ===================== Stufe: Wahrscheinlichkeit gesucht =====================
\gkopfvor{\pfstufekopf{Wahrscheinlichkeit gesucht}{\gpunkte[26 22]{\textcolor{black!80}{$\bullet$}\,\textcolor{black!25}{$\circ$}\,\textcolor{black!25}{$\circ$}\,\textcolor{black!25}{$\circ$}\,\textcolor{black!80}{$\bullet$}}}\special{pdf:dest (stufe-wahrscheinlichkeit-gesucht) [@thispage /XYZ @xpos @ypos null]}}
% kein Einstieg: keine Schülerfrage in der Gliederung (Schülerfrage, -Fälle, -Form, -Weg, -Skizze; oder Frage) – Einstieg nur von dort, nie aus der Bank (pruefung.md § 5, 12.10.)
% 2026-bb-gk-B4b-m1 (Original 2026-bb-gk-B4b, Schwere I, rang 1; herausgelöste Fassung: keine)
\begin{pfaufg}{GK~’26\\B4b}{1.}{}
In einem Supermarkt sammeln 80 \% der Kundschaft Treuepunkte. Es werden 12 Personen aus der Kundschaft zufällig ausgewählt; die Zufallsgröße $X$ gibt an, wie viele davon Treuepunkte sammeln. $X$ wird als binomialverteilt angenommen. Ermittle die Wahrscheinlichkeit dafür, dass weniger als 10 der 12 Personen Treuepunkte sammeln.
\karo{3}
\fusshilfe{\gfuss{1}{0,442}}
\end{pfaufg}
% 2022-bebb-gk-B4b-m1 (Original 2022-bebb-gk-B4b, Schwere I, rang 2; herausgelöste Fassung: keine)
\begin{pfaufg}{GK~’22\\B4b}{2.}{}
In einem Paketzentrum werden täglich mehrere Millionen Pakete umgeschlagen. Von allen Paketen gehen 12 \% zum Ziel A und 8 \% zum Ziel B, die übrigen an andere Ziele. Es werden 100 Pakete zufällig ausgewählt; $X$ ist die Anzahl der Pakete mit dem Ziel B. Bestimme mit dem Tabellenausschnitt die Wahrscheinlichkeit dafür, dass genau zehn der Pakete das Ziel B haben, und die Wahrscheinlichkeit dafür, dass weniger als zehn das Ziel B haben.
\gzwei{{\small\begin{tabular}{|l|c|c|c|c|}\hline \multicolumn{5}{|c|}{$n = 100$, $p = 0{,}08$}\\\hline $k$ & 8 & 9 & 10 & 11\\\hline $P(X \le k)$ & 0,5926 & 0,7220 & 0,8243 & 0,8972\\\hline\end{tabular}}}{\karo{3}}
\fusshilfe{\gfuss{2}{0,102; 0,722}}
\end{pfaufg}

% ===================== Stufe: n, p oder k gesucht =====================
\gkopfvor{\pfstufekopf{n, p oder k gesucht}{\gpunkte[25 23 22]{\textcolor{black!25}{$\circ$}\,\textcolor{black!80}{$\bullet$}\,\textcolor{black!25}{$\circ$}\,\textcolor{black!80}{$\bullet$}\,\textcolor{black!80}{$\bullet$}}}\special{pdf:dest (stufe-n-p-oder-k-gesucht) [@thispage /XYZ @xpos @ypos null]}}
% kein Einstieg: keine Schülerfrage in der Gliederung (Schülerfrage, -Fälle, -Form, -Weg, -Skizze; oder Frage) – Einstieg nur von dort, nie aus der Bank (pruefung.md § 5, 12.10.)
% 2023-bebb-gk-B4.1f-m1 (Original 2023-bebb-gk-B4.1f, Schwere II, rang 3; herausgelöste Fassung: keine)
\begin{pfaufg}{GK~’23\\B4.1f}{3.}{}
In einem Land arbeiten 30 \% aller Lehrkräfte an einem Gymnasium; 20 \% aller Lehrkräfte sind weiblich und arbeiten an einem Gymnasium, insgesamt sind 65 \% der Lehrkräfte weiblich. Es werden Lehrkräfte zufällig ausgewählt. Ermittle, wie viele Lehrkräfte man mindestens auswählen muss, damit mit einer Wahrscheinlichkeit von mindestens 95 \% mindestens eine davon an einem Gymnasium arbeitet.
\karo{3}
\fusshilfe{\gfuss{3}{9}}
\end{pfaufg}
% Teile eines Originals (2022-bebb-gk-B4): 2022-bebb-gk-B4d-m1, 2022-bebb-gk-B4e-m1
\begin{pfaufg}{}{4.}{}
In einem Paketzentrum werden täglich mehrere Millionen Pakete umgeschlagen. Von allen Paketen gehen 12 \% zum Ziel A und 8 \% zum Ziel B, die übrigen an andere Ziele.
% 2022-bebb-gk-B4d-m1 (Original 2022-bebb-gk-B4d, Schwere II, rang 4; herausgelöste Fassung: keine)
\begin{gbteil}{a)}{\raisebox{\dimexpr-\height+\ht\strutbox\relax}{\shortstack[r]{GK~’22\\B4d}}}{}
Ermittle, wie viele Pakete man mindestens zufällig auswählen muss, damit mit einer Wahrscheinlichkeit von mindestens 90 \% mindestens eines davon das Ziel B hat.
\karo{3}
\end{gbteil}
% 2022-bebb-gk-B4e-m1 (Original 2022-bebb-gk-B4e, Schwere II, rang 5; herausgelöste Fassung: keine)
\begin{gbteil}{b)}{\raisebox{\dimexpr-\height+\ht\strutbox\relax}{\shortstack[r]{GK~’22\\B4e}}}{}
Ein Teil der Pakete hat das Ziel C. Wählt man 20 Pakete zufällig aus, so hat mit einer Wahrscheinlichkeit von etwa 67 \% keines davon das Ziel C. Ermittle den Anteil der Pakete mit dem Ziel C an allen Paketen.
\karo{3}
\end{gbteil}
\fusshilfe{\gfuss{4}{a)~28; b)~2 \%}}
\end{pfaufg}
% 2025-bebb-gk-B4d-m1 (Original 2025-bebb-gk-B4d, Schwere II, rang 6; herausgelöste Fassung: keine)
\begin{pfaufg}{GK~’25\\B4d}{5.}{}
In einer Großstadt gelten 8,24 \% aller Haushalte als überbelegt. Es werden 800 Haushalte zufällig ausgewählt; die Zufallsgröße $X$ gibt die Anzahl der überbelegten Haushalte unter ihnen an und wird als binomialverteilt angenommen. Ermittle den kleinstmöglichen Wert von $k$, für den $P(X \le k)$ mehr als 90 \% beträgt.
\karo{3}
\fusshilfe{\gfuss{5}{76}}
\end{pfaufg}

\pfblattende{10}{26 & 25 & 24 & 23 & 22 & 21 & 20 & 19 & 18 & 17}{B4b & B4d &  & B4.1f & B4b B4d B4e & B4a B4e & B4.1c B4.2a & B4.1a B4.1b & B3.2a B3.2d & }%
{\pfausschnitt{Bernoulli-Kette}{Binomialwahrscheinlichkeit}{Säulendiagramm und Erwartungswert}}
\end{document}